\documentclass[11pt]{article}

\usepackage[T1]{fontenc}
\usepackage[utf8]{inputenc}
\usepackage{amsmath,amssymb,amsthm}
\usepackage{booktabs}
\usepackage{enumitem}
\usepackage[margin=1in]{geometry}
\usepackage{microtype}
\usepackage{xcolor}
\usepackage{hyperref}

% The offline Tectonic bundle lacks the 7pt Computer Modern symbol outline.
% Use the bundled 8pt outlines for script mathematics at the nearby text sizes.
\DeclareMathSizes{5}{8}{8}{8}
\DeclareMathSizes{6}{8}{8}{8}
\DeclareMathSizes{7}{8}{8}{8}
\DeclareMathSizes{8}{8}{8}{8}
\DeclareMathSizes{9}{9}{8}{8}
\DeclareMathSizes{10}{10}{8}{8}
\DeclareMathSizes{11}{11}{8}{8}
\DeclareMathSizes{12}{12}{8}{8}
\DeclareFontShape{OML}{cmm}{m}{it}{<5-9>cmmi8<9->cmmi10}{}
\DeclareFontShape{OMS}{cmsy}{m}{n}{<5-9>cmsy8<9->cmsy10}{}
\DeclareFontShape{OT1}{cmr}{m}{n}{<5-9>cmr8<9->cmr10}{}
\DeclareFontShape{OT1}{cmtt}{m}{n}{<5-9>cmr8<9->cmr10}{}

\hypersetup{
  colorlinks=true,
  linkcolor=blue!55!black,
  citecolor=blue!55!black,
  urlcolor=blue!55!black,
  pdftitle={Posterior Amplification in CS-128: Equivalent-Key Recovery and a Fresh-Message Forgery from 2.3 Million Signatures},
  pdfauthor={Martin Feussner}
}

\newtheorem{theorem}{Theorem}
\newtheorem{lemma}{Lemma}
\newtheorem{corollary}{Corollary}
\newtheorem{remark}{Remark}
\newcommand{\R}{\mathcal R}
\newcommand{\F}{\mathcal F}
\newcommand{\HB}{\operatorname{HighBits}}
\newcommand{\LB}{\operatorname{LowBits}}
\newcommand{\E}{\mathbb E}
\newcommand{\Prb}{\mathop{\rm Pr}}
\newcommand{\CS}{\mathsf{CS}}
\newcommand{\Sign}{\mathsf{Sign}}
\newcommand{\Verify}{\mathsf{Verify}}

\title{\bfseries Posterior Amplification in CS-128:\\
Equivalent-Key Recovery and a Fresh-Message Forgery\\[0.35em]
\large 2.3 Million Reference Signatures on One Fixed Key}
\author{Martin Feussner\\
Selmer Center, University of Bergen\\
\texttt{martin.feussner@uib.no}}
\date{28 September 2026}

\begin{document}
\maketitle

\noindent\textbf{Provenance and review status.}
This work was implemented and tested during an AI-run audit by OpenAI Codex using
Daybreak Blue at maximum reasoning effort.  The public report by Yijian Liu on behalf of
Xianhui Lu established priority for the underlying CS compressed-signature leakage and its
claimed key-recovery consequence~\cite{liu-lu,ngcc-report}.  Their attack code and transcripts
were not available.  This note gives a self-contained derivation and independent
implementation, adds a verifier-conditioned approximate-posterior score and a deterministic
public completion procedure, and records a concrete fresh-message forgery from the recovered
equivalent key.  It does not claim priority for the general key-recovery consequence of the
leakage.  The derivation, implementation, experiment, and manuscript were produced by the
AI.  At the date above, no human cryptanalyst has independently verified the argument,
implementation, measurements, or novelty assessment.  This is a preliminary disclosure for
human review and reproduction.

\begin{abstract}
CS is a Fiat--Shamir-with-aborts lattice signature whose full construction compresses the
public product and omits the final response vector from the signature.  Verification publicly
reconstructs a noisy vector
\[
 z'_2=y_2-c_r(b_0-s_2)+\nu,
\]
where $b_0$ and $s_2$ belong to the signing key and $c_r$ contains six hidden bimodal signs.
The accepted $z_0=y_0+c_r$ is statistically correlated with those signs.  Liu and Lu used
this fact to report CS-128 key recovery from $2^{27}$ signatures.

We give an independent implementation and refinement using
$2{,}300{,}000<2^{21.14}$ signatures, about 58 times fewer, on one fixed CS-128 key.  One
million signatures from the unmodified reference signer give a
rough public estimate of $d=b_0-s_2$.  For every later transcript we enumerate the $2^6$ possible hidden sign
vectors and combine the unconditioned $D_{5,5}$ mask likelihood from $z_0$ with the residual energy
$\lVert z'_2+c_r\widehat d\rVert^2$ in an approximate posterior score.  Regression of public
$z_1$ on the resulting approximate-posterior mean over 1.3 million further signatures
gives a ternary candidate differing from the originally sampled $s_1$ at two coefficients.
A deterministic public completion ladder changes one coefficient and obtains a compatible
equivalent signing key.

The completed candidate $u$ determines compatible values $s_2^*$ and $b_0^*$.  We serialized
the resulting equivalent signing key, signed a message absent from the training set, and obtained
acceptance from the reference CS-128 verifier.  The experiment ran on a four-core Intel
Core i5-6500; Section~\ref{sec:experiment} reports its measured time.  The attack follows the
normative signing and verification equations and demonstrates an end-to-end attack through the
specified design on the tested key.  Success probability over independently
generated keys has not been measured.
\end{abstract}

\section{Introduction}

CS (``Compact-Simple''), a signature candidate in the Next-generation Commercial
Cryptographic Algorithms Program (NGCC), is a module-lattice construction with bimodal
rejection sampling and two compression mechanisms~\cite{cs-spec}.  The public key contains
only the high part $b_1$ of the module-LWE product
\[
 b=A's_1+s_2=2^\beta b_1+b_0,
\]
while a signature transmits $z_0$, $z_1$, a hint, and a Fiat--Shamir seed.  The verifier
reconstructs the missing $z_2$ component from those values.  This saves public-key and
signature bandwidth, but the reconstruction retains a product between secret low information
and the hidden signs chosen by the iterative bimodal sampler.

Liu and Lu publicly identified this leakage and reported complete CS-128 recovery from
$2^{27}$ signatures generated in 3.25 hours on 16 cores~\cite{liu-lu}.  The public NGCC
tracker verified the algebraic identity but classified the statistical recovery as a lead
because no code or transcript was supplied~\cite{ngcc-report}.  Their report estimates the
hidden bimodal aggregate from $(z_0,c)$, regresses the reconstruction to obtain low secret
information, and then uses a cross moment to obtain $s_1$.

This note independently implements and strengthens that attack.  Its main observation is that
the public reconstruction $z'_2$ should be used twice: first as the regression target for a
rough low-part estimate and then as side information when inferring each transcript's hidden
signs.  CS-128 has only
\[
 N=\left\lceil\frac{\tau}{\tau'}\right\rceil
  =\left\lceil\frac{23}{4}\right\rceil=6
\]
hidden block signs.  Exact enumeration therefore costs only 64 likelihood evaluations per
signature.  Even a rough estimate of the low secret term gives a strong 768-coordinate energy
test for those 64 choices.  The resulting approximate-posterior score is informative enough
that regression on $z_1$, followed by the deterministic public completion ladder, constructs
an equivalent signing key in the experiment below.

The candidate-producing path and its stopping decision do not read the target secret.  A
candidate $s_1$ is accepted only if every coefficient of the public rounded product can be
completed by a ternary $s_2$, after which the equivalent serialized key must produce a
reference-verifier-accepted signature on a fresh message.

The attack has the following concrete properties.
\begin{enumerate}[leftmargin=1.7em]
  \item Every estimator input is reconstructed from the public key and ordinary signatures.
  \item The sample complexity is $2{,}300{,}000<2^{21.14}$ signatures, about $58.4$ times
        below the previously reported $2^{27}$ experiment.
  \item Statistics are accumulated in streaming form; the largest per-worker state file is
        110,608 bytes, and no signature corpus is retained.
  \item The final completion check and the equivalent-key construction use only the public key.
  \item The final result is an accepted fresh-message forgery, rather than a coefficient-match
        diagnostic alone.
\end{enumerate}

Section~\ref{sec:scheme} isolates the relevant CS equations.
Section~\ref{sec:identity} proves the public leakage identity.
Sections~\ref{sec:baseline} and~\ref{sec:amplification} give the two-stage estimator.
Section~\ref{sec:completion} describes public correction and equivalent-key construction.
Section~\ref{sec:experiment} reports the full-parameter experiment and conformance checks.
Sections~\ref{sec:scope} and~\ref{sec:repair} discuss scope and repair.

\section{CS-128 notation and signing equations}
\label{sec:scheme}

Let
\[
  \R_q=\mathbb Z_q[x]/(x^n+1),
  \qquad (n,q,k,\ell)=(256,32257,3,3).
\]
Products below are negacyclic ring products, componentwise across module coordinates where
appropriate.  We use centered representatives whenever a small integer value is required.
The remaining relevant CS-128 parameters are
\begin{equation}
  (\tau,\tau',N)=(23,4,6),\qquad 2^\beta=32,\qquad \alpha=2016. \label{eq:params}
\end{equation}
The secret vectors $s_1\in\{-1,0,1\}^{\ell\times n}$ and
$s_2\in\{-1,0,1\}^{k\times n}$ satisfy
\begin{equation}
  b=A's_1+s_2=2^\beta b_1+b_0\pmod q,                  \label{eq:key}
\end{equation}
where the public key is $(\rho_0,b_1)$, $A'=\mathsf{ExpandA}(\rho_0)$, and $b_0$ is
stored only in the secret key.  Define the effective hidden low term
\begin{equation}
  d=b_0-s_2=A's_1-2^\beta b_1\pmod q.                 \label{eq:d}
\end{equation}
Although Equation~\eqref{eq:d} becomes public once $s_1$ is known, $d$ is secret during
the statistical phase.

For one signing attempt, the signer samples independent mask vectors
\[
 y_0\leftarrow D_{5,5},\qquad
 y_1\leftarrow D_{10,9}^{\ell},\qquad
 y_2\leftarrow D_{6,9}^{k}.
\]
Here a coefficient of $D_{a,b}$ is the difference of two sums of $a$ independent uniform
$b$-bit integers.  The Fiat--Shamir challenge $c$ has $\tau=23$ signed nonzero
coefficients.  The sampler processes the challenge support in six consecutive blocks
$S_0,\ldots,S_5$ of sizes $4,4,4,4,4,3$.  It chooses a hidden sign
$\varepsilon_g\in\{+1,-1\}$ for each block, subject to the specified rejection test, and
forms
\begin{equation}
  r=c_r=\sum_{g=0}^{5}\varepsilon_g c\vert_{S_g}.       \label{eq:r}
\end{equation}
Thus $r$ has the public support and coefficient magnitudes of $c$, while six signs are hidden.
An accepted response satisfies
\begin{equation}
 z_0=y_0+r,\qquad z_1=y_1+rs_1,\qquad z_2=y_2+rs_2.     \label{eq:responses}
\end{equation}
Only $z_0$ and $z_1$ are encoded.  The hint lets the verifier reproduce the high part of the
commitment $w$ and reconstruct a bounded vector denoted $z'_2$ below.  All quantities needed
to compute $z'_2$ are public in a verification transcript.

\section{The verifier-visible leakage identity}
\label{sec:identity}

We first restate the correctness calculation in a form useful for estimation.  Let
$\operatorname{LowBits}_h(w,\alpha)$ denote the centered low part associated with the hint
decomposition, and let $\operatorname{LSB}(w)$ denote the parity lift used by CS.  Define
\begin{equation}
  \nu=\frac{-\operatorname{LowBits}_h(w,\alpha)+\operatorname{LSB}(w)}{2}. \label{eq:nu}
\end{equation}
The parity conventions make the numerator even coefficientwise.

\begin{lemma}[Public reconstruction]
For every honestly generated accepted signature, Algorithm~12 publicly reconstructs
\begin{equation}
 z'_2=z_2-rb_0+\nu=y_2-rd+\nu\pmod{\pm q}.             \label{eq:leak}
\end{equation}
\end{lemma}

\begin{proof}
The signing hint is the difference between the high parts of $w$ and
$w-2z_2+2rb_0$.  In verification, substitution of
$z_0=y_0+r$ and $z_1=y_1+rs_1$ into the public linear expression gives
\[
 v=w-2z_2+2rb_0\pmod{2q}.
\]
The term involving the original challenge has the same parity as $r$, because changing a
challenge sign changes a coefficient by two; it therefore vanishes after multiplication by
$q$ modulo $2q$.  The hint recovers the high part of $w$, and the parity check recovers its
least significant bit.  Algorithm~12 consequently computes
\[
 \frac{\alpha\HB_h(w)+\operatorname{LSB}(w)-v}{2}
 =z_2-rb_0+
   \frac{-\operatorname{LowBits}_h(w,\alpha)+\operatorname{LSB}(w)}{2}.
\]
Using $z_2=y_2+rs_2$ and $d=b_0-s_2$ proves Equation~\eqref{eq:leak}.
\end{proof}

Equation~\eqref{eq:leak} is the design-level leakage channel.  The verifier needs $z'_2$ for
its norm check, so the attacker obtains a noisy linear observation of $rd$.  The sign vector
$r$ is hidden, but it is statistically visible through $z_0=y_0+r$.  The first stage obtains a
rough $d$ using only that weak visibility.  The second stage feeds the rough $d$ back into
Equation~\eqref{eq:leak} to infer $r$ much more accurately.

The scale used by the second-stage likelihood is also explained by public parameters.  A
$D_{6,9}$ coefficient has variance
\begin{equation}
 \operatorname{Var}(D_{6,9})
 =\frac{6(2^{18}-1)}{6}=262143.                         \label{eq:y2var}
\end{equation}
Treating the centered hint low part as approximately uniform contributes about
$\alpha^2/48=84672$ after division by two.  Their sum is about $346815$.  We therefore fix
the public likelihood parameter $\sigma^2=350000$ from public parameters; no fitted secret
scale is needed.

\section{Stage one: a rough low-term estimate}
\label{sec:baseline}

Let $p_0(a)=\Prb[Y_0=a]$ be the exact one-coefficient probability mass function of
$D_{5,5}$.  It is computed once by integer convolution of five uniform 5-bit variables with
five negated copies.  For a transcript $t$ and one challenge block $S_g$, define
\begin{align}
 \lambda_{t,g}(+1)&=\sum_{i\in S_g}\log p_0(z_{0,t,i}-c_{t,i}),\\
 \lambda_{t,g}(-1)&=\sum_{i\in S_g}\log p_0(z_{0,t,i}+c_{t,i}).    \label{eq:blockll}
\end{align}
Ignoring the final joint rejection event gives the public soft sign
\begin{equation}
 m_{t,g}=\tanh\!\left(\frac{\lambda_{t,g}(+1)-
                                  \lambda_{t,g}(-1)}{2}\right),    \label{eq:softsign}
\end{equation}
and a ring element $\widehat r_t$ whose coefficient at $i\in S_g$ is
$m_{t,g}c_{t,i}$ and whose other coefficients are zero.

We diagonalize negacyclic multiplication with the twisted Fourier transform.  For each
frequency $f$ and public reconstruction component $a$, centered least squares gives
\begin{equation}
 \widehat d_{a,f}=-
 \frac{\sum_t \overline{(\widehat r_{t,f}-\overline r_f)}
                    (z'_{2,t,a,f}-\overline z'_{2,a,f})}
      {\sum_t|\widehat r_{t,f}-\overline r_f|^2}.                   \label{eq:dreg}
\end{equation}
An inverse twisted transform returns coefficient estimates, which we round to the small
centered range allowed by Equation~\eqref{eq:key}.  This estimate is deliberately crude.  In
our one-million-signature reference pass it had mean absolute error $1.665$, root mean square
error $2.080$, and maximum error $6.170$ against the diagnostic value of $d$; only 155 of 768
coefficients rounded exactly.  Its purpose is to score hidden sign vectors, where the energy
is accumulated over all 768 reconstruction coordinates.

The refinement score depends on this completed calibration estimate.  Our streaming
implementation therefore discards the calibration transcripts and applies the score to a
disjoint second pass.  Reusing the first million signatures would instead require retaining
those transcripts until calibration finishes.

\section{Stage two: verifier-conditioned approximate-posterior amplification}
\label{sec:amplification}

For each new signature, enumerate
$\varepsilon=(\varepsilon_0,\ldots,\varepsilon_5)\in\{\pm1\}^6$ and construct
$r_\varepsilon$ as in Equation~\eqref{eq:r}.  Assign the log weight
\begin{equation}
 \log W_t(\varepsilon)=
   \sum_{g=0}^{5}\lambda_{t,g}(\varepsilon_g)
   -\frac{\left\lVert z'_{2,t}+r_\varepsilon\widehat d\right\rVert_2^2}
          {2\sigma^2},
 \qquad \sigma^2=350000.                                           \label{eq:weight}
\end{equation}
The first term is the exact unconditioned $D_{5,5}$ mask likelihood used in stage one.  For the correct
sign vector and an accurate $\widehat d$, Equation~\eqref{eq:leak} makes the second term the
energy of $y_2+\nu$; a wrong sign vector adds a structured convolution by $d$.  Because this
likelihood intentionally omits rejection-induced dependence, $W_t$ is an approximate posterior
score rather than an exact posterior probability.  Its normalized mean
\begin{equation}
 \widetilde r_t=
 \frac{\sum_{\varepsilon}W_t(\varepsilon)r_\varepsilon}
      {\sum_{\varepsilon}W_t(\varepsilon)}                         \label{eq:rpost}
\end{equation}
is computed by a log-sum-exp normalization.  Its inputs are the public transcript, public
parameters, and the stage-one estimate.

Now use the other public response in Equation~\eqref{eq:responses}.  At each frequency $f$
and secret-vector component $j$, centered least squares gives
\begin{equation}
 \widehat s_{1,j,f}=
 \frac{\sum_t\overline{(\widetilde r_{t,f}-\overline{\widetilde r}_f)}
                   (z_{1,t,j,f}-\overline z_{1,j,f})}
      {\sum_t|\widetilde r_{t,f}-\overline{\widetilde r}_f|^2}.    \label{eq:sreg}
\end{equation}
After the inverse twisted transform, every coefficient is rounded and clipped to
$\{-1,0,1\}$.

The likelihood in Equation~\eqref{eq:weight} is intentionally simple: it treats the residual
as spherical and does not model rejection-induced dependence.  The experiment therefore
evaluates the approximate-posterior score through the resulting public-key recovery and forgery,
without treating a hidden-sign diagnostic as part of the attack evidence.

The status of the argument is consequently mixed but explicit.  The public-reconstruction
identity in Equation~\eqref{eq:leak}, the completion predicate, and the equivalent-key lemma
are algebraic consequences of the specified algorithms.  The spherical likelihood and
$\sigma^2=350000$ are modelling choices.  The public ladder's first successful sample count is
an empirical result on one fixed key, rather than a concentration bound or a theorem over the
key-generation distribution.

\section{Public correction and equivalent-key construction}
\label{sec:completion}

Direct rounding need not be perfect.  The public key supplies an exact and inexpensive
completion predicate.  For a ternary candidate $u\in\{-1,0,1\}^{\ell\times n}$, define
\begin{equation}
 I(u)=\#\left\{(a,i):
   \nexists e\in\{-1,0,1\}\ \text{such that}\
   \HB_{2^\beta}\bigl((A'u)_{a,i}+e\bigr)=(b_1)_{a,i}
 \right\}.                                                        \label{eq:invalid}
\end{equation}
This uses the exact submitted rounding rule, including boundary representatives.
Every valid signing key has $I(s_1)=0$.  Whenever $I(u)=0$, one admissible $e$ at every
coordinate forms a ternary vector $s_2^*$ and
\begin{equation}
  b_0^*=A'u+s_2^*-2^\beta b_1                                  \label{eq:equiv}
\end{equation}
forms its matching low part.  Thus a zero of $I$ is already an equivalent signing key; it
does not have to equal the originally sampled key.

The implementation contains a deterministic public fallback ladder for runs in which direct
rounding does not give $I(u)=0$.  It first tests scales from $0.7$ through $1.3$ in increments
of $0.0002$, rounding and clipping after each scale.  It then ranks coefficients by the increase
in squared error to the closest alternative ternary value and tests: both alternatives at each
of the first 64 positions; all two-position alternatives among those 64 positions; and all
three-position alternatives among the first 24 positions.  The coefficient stage contains at
most
\[
 2\cdot64+4\binom{64}{2}+8\binom{24}{3}=24{,}384
\]
public candidates.  Every decision uses only the continuous public estimate and $I(u)$.  This
is also the candidate ladder used by the 100,000-signature stopping rule in
Section~\ref{sec:experiment}; the stopping decision is the public condition $I(u)=0$ followed
by a verifier-accepted fresh-message signature.  Secret-coefficient diagnostics do not choose
the stopping point or a candidate.

Before serialization, $b_0^*$ is formed from the standard representative of
$A'u+s_2^*\bmod q$ using the submitted implementation's exact centered/\texttt{LowBits}
representative convention.  The code applies \texttt{Power2RoundVector}, then checks its high
part coefficientwise against $b_1$.  To serialize an equivalent secret key, choose an
arbitrary fresh secret seed $K$, copy the public matrix seed $\rho_0$, compute
$\mathit{tr}=H(\mathit{pk})$, and encode
$(u,s_2^*,K,\mathit{tr},\rho_0,b_0^*)$.  The submitted byte format additionally stores the
public high part $b_1$ (named \texttt{t1} in the code), which is copied from the public key.
The reference signer can then sign arbitrary messages for the original public key.

\begin{lemma}[Equivalent-key completion]
If a ternary candidate $u$ satisfies $I(u)=0$, the public completion procedure constructs a
well-formed CS-128 signing key for the target public key.
\end{lemma}

\begin{proof}
Equation~\eqref{eq:invalid} supplies $s_2^*$ coefficientwise and
Equation~\eqref{eq:equiv} supplies the corresponding low part.  These values satisfy the same
public rounding relation as the target key.  The candidate $u$ and the completed $s_2^*$ have
the required ternary form.  The public seed $\rho_0$, public high part $b_1$, and
$H(\mathit{pk})$ are copied, while the signing-randomness seed $K$ is freely chosen.
Encoding these values with $b_0^*$ therefore
gives a signing key satisfying the CS-128 key relation.  Whenever the signing algorithm returns
a response, the specified correctness calculation consequently verifies it under the target
public key.
\end{proof}

In the experiment below, the deterministic public candidate ladder produced a vector with
$I(u)=0$.  The resulting key signed a fixed message absent from every signing-oracle query,
and Algorithm~12 accepted the signature under the target public key.

The completion lemma concerns the specified algebraic key relation.  Its byte-level realization
uses the submitted \texttt{Power2RoundVector} and \texttt{skEncode} routines, then checks the
result with \texttt{CS\_Verify}.  The experiment checks both layers.

\section{Full-parameter experiment}
\label{sec:experiment}

\subsection{Setup}

The experiment used the submitted CS-128 parameter set and reference source.  One fixed,
deterministically generated test key was shared by all worker streams.  Messages were distinct
16-byte counter-and-worker strings.  Each worker accumulated Fourier moments independently;
the small state files were merged without retaining signatures.  Ground-truth values were recorded only to
report estimator accuracy.  They were not inputs to the approximate posterior scores,
regressions, candidate ordering, public completion predicate, equivalent-key construction, or
verifier decision.

For exact reproduction, key-generation seed byte $i$ and worker-$w$ seed byte $i$ were
respectively
\[
 \mathtt{0x51}+7i\pmod {256},\qquad
 \mathtt{0xa7}+11i+37w\pmod {256}.
\]
The x86-64 host encoded integers in little-endian order.  The first eight message bytes held
query index $j$; the final eight held
\[
j\mathbin{\mathsf{xor}}\mathtt{0x43532d313238}\mathbin{\mathsf{xor}}(w\ll48).
\]
The worker sets $w\in\{1,2,3,4,101,102,103\}$ are disjoint, and multiplication
by 37 permutes bytes modulo 256, so all seven DRNG seeds differ.

The forgery target was the 42-byte ASCII message
\texttt{CS-128 recovered-key fresh-message forgery}.  Its length alone distinguishes it from
every 16-byte signing query, so freshness does not depend on a probabilistic collision argument.

The calibration pass generated 250,000 signatures per worker through the submitted reference
$\CS.\Sign$ path, using worker streams 1--4.  Refinement used the disjoint streams 101, 102,
and 103 in that fixed order.  Streams 101 and 102 contributed 500,000 signatures each; stream
103 stopped after 300,000 because the attack succeeded.  A cumulative state was frozen every 100,000 signatures.  At each
point the deterministic public candidate ladder was run; the attack transcript ends at its first
verifier-accepted fresh-message signature.  Section~\ref{sec:experiment} reports that checkpoint
and its immediately preceding failure.  Secret diagnostics do not select the endpoint.  Every signature
in those states was produced by the same unmodified reference signer and sampler, and no message
or randomness stream was reused between the two passes.  The collector linked copies
of the submitted \texttt{cs.c} and \texttt{sampling.c} files directly.  Those copies differ
from the official archive only in CRLF line endings; their line-ending-normalized SHA-256
hashes match the official files and are recorded below.  No alternative sampler was used, and
verifier acceptance is not used as evidence of distributional identity.  The forged signature
was likewise produced by the submitted reference signer and checked by its reference verifier.

For transparency, the collector has no sampler instrumentation, hidden-sign trace, or fast
mode.  Its refinement stage refuses to run without the public stage-one estimate.  A separate
deterministic algebra test compares the optimized Gram-matrix evaluation of all 64 residual
energies with their direct 768-coordinate evaluation; all tested scores agree exactly.  In a
negative-control build, every ground-truth array supplied to the reporting routine is replaced
by unrelated fixed values immediately after the public states are merged.  The first successful
checkpoint, correction index, completed candidate, and accepted fresh-message forgery are
unchanged.  This control confirms that
the diagnostic comparisons do not drive recovery or completion.

\subsection{Results}

Table~\ref{tab:results} summarizes the attack on a four-core Intel Core i5-6500 desktop with
16~GB of memory, Linux 7.0.0, GCC 13.3.0, and
\texttt{-O3 -flto -fcommon -std=gnu11}.  The managed execution environment scheduled the
workers below the machine's full four-core capacity, so we report both aggregate user CPU and
observed elapsed time.  The final worker's elapsed time includes deliberate pauses at the
100,000-signature gates while the public tests ran.  The public-ladder row is a complete replay
of all 13 tested refinement prefixes, including equivalent-key encoding, fresh signing, and
reference verification at the successful prefix.

\begin{table}[ht]
\centering
\caption{CS-128 full-parameter attack results.}
\label{tab:results}
\small
\begin{tabular}{@{}lrrrl@{}}
\toprule
Phase & Signatures & User CPU & Elapsed & Principal result \\
\midrule
Reference calibration & $1{,}000{,}000$ & $28.18$ min & $35.15$ min &
  $\widehat d$ MAE $1.665$, RMSE $2.080$ \\
Approx. posterior refinement & $1{,}300{,}000$ & $37.50$ min & $74.53$ min &
  public transcript processing \\
Public ladder replay & $0$ & $0.12$ min & $0.25$ min &
  depth one; $I(u)=0$; verifier accepts \\
\bottomrule
\end{tabular}
\end{table}

The final diagnostic record was
\begin{verbatim}
approx_posterior_s1_errors=2/768
approx_posterior_s1_maxerr=0.525291
public_correction_depth=1 index=287
public_key_consistency_invalid=0/768
equivalent_key_fresh_forgery=ACCEPT
\end{verbatim}
Direct ternary rounding differed from the originally sampled $s_1$ at two of 768 coefficients.
This comparison is diagnostic only.  The successful path used the raw regression values,
ternary rounding, and the fixed public predicate in Equation~\eqref{eq:invalid}.  The scale
sub-ladder found no valid key; the coefficient sub-ladder changed index 287 and obtained
$I(u)=0$ at depth one.  It did not fit or apply a secret-derived scale, and the resulting
equivalent $s_1$ need not equal the originally sampled one.

The stage-one workers each retained 79,888 bytes and the refinement workers each retained
110,608 bytes.  Consequently, the attack can consume signatures online and discard them
immediately.  Its total sample count is $2{,}300{,}000$, or about $2^{21.13}$, compared with
$2^{27}$ in the earlier report.  Signature collection consumed 3,940.53 seconds of aggregate
user CPU and 6,580.46 seconds (109.67 minutes) of observed staged elapsed time, including the
checkpoint pauses.

The frozen public ladder tested refinement prefixes in increments of 100,000.  It rejected
every prefix from 100,000 through 1,200,000 and first accepted at 1,300,000.  Thus 2.3 million
total signatures are the first success at this resolution for the fixed key, calibration budget,
estimator, and candidate ladder.  This is not a general minimality result: the one-million
calibration budget was not varied, intermediate counts were not tested, and another public
completion algorithm might succeed from an earlier prefix.

\subsection{Reproducibility record}

The principal artifact hashes are listed in Table~\ref{tab:hashes}.  The evidence manifest
fixes every streaming checkpoint and timing record; the ladder and result files record the
public stopping decision, completion, and accepted forge.

\begin{table}[ht]
\centering
\caption{SHA-256 hashes of the attack source and final evidence.}
\label{tab:hashes}
\small
\begin{tabular}{@{}p{0.25\textwidth}p{0.68\textwidth}@{}}
\toprule
Artifact & SHA-256 \\
\midrule
\texttt{leak\_reference.c} &
\texttt{dcafca8a30fe43f656dcd71b701b9c61}\newline
\texttt{78a8dffaf1d75501e6e2a22926a4ea3c} \\
Official \texttt{CS.zip} archive &
\texttt{c790d31cd4a288990f3d692381ed721a}\newline
\texttt{06641938a02dfc2e0435b7d323475cef} \\
Submitted \texttt{cs.c}\newline and \texttt{sampling.c}\newline
(LF-normalized) &
\texttt{ac771cc2f138158680d1504f92bf2398}\newline
\texttt{c3f8fc85b9a6137c8cebcaa18bf96ff0}\newline
\texttt{bfd6abad31099e2a045644d500d3d1df}\newline
\texttt{4f90ae98d50a5c10ebf0a79b0416f541} \\
\texttt{allref\_coarse.raw} &
\texttt{45e24ce4106e3c32772674dd7ad3d7ff}\newline
\texttt{58701fe3dfdc55a8de86cc9b5b15fb15} \\
\texttt{evidence/SHA256SUMS} &
\texttt{91ce157728b2c5b794c7d3e06768d624}\newline
\texttt{d9273566d302e3de2dab6ec3318ebce4} \\
\texttt{public\_ladder.txt} &
\texttt{ea86e153f5a631529470142317fa212b}\newline
\texttt{5236bcbbc78a2d05489f741682a10131} \\
\texttt{cs128\_public\_ladder\_}\newline\texttt{result.txt} &
\texttt{f8527834c8fe0d4aac07a687f6657a37}\newline
\texttt{cf6c03928b4c09181a95af56ceeba28d} \\
\bottomrule
\end{tabular}
\end{table}

The frozen states reproduce rejection at all twelve earlier 100,000-signature refinement
checkpoints, the depth-one public correction at 1.3 million refinement signatures, zero
public-key violations, and the accepted forgery when merged by the listed source.

\section{Security impact and scope}
\label{sec:scope}

For the tested key, the result is a complete specification-conforming many-query attack: it
recovers an equivalent signing key and produces a valid signature on a fresh message.  It uses
the normative public-key compression, iterative bimodal sampler, hint equation, and verifier
reconstruction in Algorithms~10--12.  The leakage identity is therefore a design property;
repairing rANS canonicality, signature-length handling, memory safety, or another
implementation defect does not remove Equation~\eqref{eq:leak}.

The final key is equivalent rather than necessarily byte-identical to the original key.  This
distinction has no protective value: it signs arbitrary messages under the victim public key.
The transcript messages can be chosen by the attacker, and the final message is fresh.  The
experiment therefore satisfies the forgery condition for this fixed public key.  Because only
one generated key was tested, it does not estimate the probability of success in the full
EUF-CMA experiment over key generation.

This experiment evaluates CS-128 only.  CS-256 and CS-512 have different ring dimensions,
block counts, mask parameters, and compression parameters.  The same algebraic identity and
approximate-posterior construction apply, but this note makes no laptop sample-complexity claim for those
sets.  The result also does not depend on the separate challenge-sign-blind search attack
recorded for CS~\cite{ngcc-report}.

The attack is statistical and its success probability across independently generated keys has
not been measured.  The experiment uses one fixed key and deterministic streams, so it does
not measure key-averaged reliability.  It is nevertheless a complete submitted-parameter
experiment on that key: it uses public transcript processing, public-key-only completion, and
an accepted fresh-message forgery.

\section{Repair considerations}
\label{sec:repair}

The failure comes from combining three facts: $r$ is only weakly hidden by a narrow mask in
$z_0$, the public-key low part is secret, and verification exposes a noisy linear function of
$r$ times that secret low information.  Increasing a norm bound or changing the byte encoding
does not alter this channel.

A repair needs a new construction and security analysis.  At minimum, the verification
transcript must cease to reveal a usable relation of the form
\[
 \text{public reconstruction}=\text{independent noise}-r\cdot
 \text{secret low term}.
\]
Possible design directions include publishing enough of the rounded relation that the omitted
term is public, transmitting and checking a response component whose masking distribution is
proved independent of the key, or redesigning the hint and compression together.  Each choice
changes sizes, correctness margins, or underlying assumptions.  Reparameterizing the current
estimator alone is not a repair: the attack exploits the scheme's actual joint transcript
distribution.

\section{Conclusion}

CS-128's compressed verifier exposes a noisy product between six hidden bimodal signs and
secret low key information.  A one-million-signature public regression gives a rough model of
that term.  Enumerating only 64 possibilities per later transcript then turns the verifier's
own reconstructed vector into a strong approximate-posterior score for the hidden signs.  After
1.3 million further signatures, the deterministic public ladder changes one rounded coefficient,
constructs an equivalent key, and produces a reference-verifier-accepted fresh-message
signature.  The complete experiment uses 2.3 million signatures.  On the four-core i5-6500,
collection consumed 65.68 minutes of aggregate user CPU and 109.67 minutes of observed staged
elapsed time, including deliberate checkpoint pauses in the managed execution environment.

\begin{thebibliography}{9}

\bibitem{cs-spec}
CS submission team.
\newblock \emph{CS: Compact-Simple Signature Scheme}, NGCC submission specification,
Algorithms 10--12 and Table 3, 2026.
\newblock Available with the official artifact mirror at
\url{https://github.com/ngcc-dev/ngcc-harness/blob/main/sign-07/sign-07-spec.pdf}.

\bibitem{liu-lu}
Yijian Liu, on behalf of Xianhui Lu.
\newblock ``Leakage issue in the CS signature scheme,'' NGCC PKC Forum post,
24 September 2026.  AI assistance disclosed by the authors.
\newblock Linked from the public CS report at
\url{https://ngcc.dev/reports/sign-07.html}.

\bibitem{ngcc-report}
NGCC artifact project.
\newblock ``CS (sign-07): compressed signatures may leak the CS-128 signing key,''
public report and reproduction status, 2026.
\newblock \url{https://ngcc.dev/reports/sign-07.html}.

\end{thebibliography}

\end{document}
